\cleardoublepage
\section{势 和 场}
\subsection{势表述}
\subsubsection{标势与矢势}
\begin{align*}
    \boldsymbol{B}&=\boldsymbol{\nabla}\times\boldsymbol{A}\\
\boldsymbol{\nabla}\times\boldsymbol{E}&=-\frac{\partial}{\partial t}\boldsymbol{\nabla}\times\boldsymbol{A}\\
\boldsymbol{\nabla}\times\left(\boldsymbol{E}+\frac{\partial\boldsymbol{A}}{\partial t}\right)&=0
\end{align*}
括号中的量，与单独的$\boldsymbol{E}$不同，它可写成一个标量的梯度：
\begin{align}
    \boldsymbol{E}+\frac{\partial\boldsymbol{A}}{\partial t}&=-\boldsymbol{\nabla}V\\
    \boldsymbol{E}&=-\boldsymbol{\nabla}V-\frac{\partial\boldsymbol{A}}{\partial t}\\
    \boldsymbol{\nabla}\cdot\boldsymbol{E}&=-\nabla^2V-\frac{\partial\boldsymbol{\boldsymbol{\nabla}\cdot A}}{\partial t}\\
    \frac{1}{\varepsilon_0}\rho&=-\nabla^2V-\frac{\partial\boldsymbol{\boldsymbol{\nabla}\cdot A}}{\partial t}\\
    \boldsymbol{\nabla}\times\left(\boldsymbol{\nabla}\times\boldsymbol{A}\right)&=\mu_0\boldsymbol{J}
    -\mu_0\varepsilon_0\boldsymbol{\nabla}\frac{\partial V}{\partial t}
    -\mu_0\varepsilon_0\frac{\partial^2\boldsymbol{A}}{\partial t^2}
\end{align}
\subsubsection{规范变换}
只要不影响$\boldsymbol{E}$和$\boldsymbol{B}$，我们可以自由地对$V$和$\boldsymbol{A}$附加额外条件。假设我们有两套势$(V,\boldsymbol{A})$和$(V',\boldsymbol{A}')$，它们对应相同的电场和磁场。设
$$\boldsymbol{A}'=\boldsymbol{A}+\alpha,\, V'=V+\beta$$
因为$\boldsymbol{A}$与$\boldsymbol{A}'$要给出同样的$\boldsymbol{B}$，它们的旋度必须相等，故
$$\boldsymbol{\nabla}\times\boldsymbol{\alpha}=0$$
因此我们可以把$\boldsymbol{\alpha}$写成某个标量的梯度
$$\boldsymbol{\alpha}=\boldsymbol{\nabla}\lambda$$
这两个势也要给出同样的$\boldsymbol{E}$，所以
$$\boldsymbol{\nabla}\left(\beta+\frac{\partial\lambda}{\partial t}\right)=0$$
括号中的项因此不依赖于坐标（然而它可以依赖于时间），称它为$k(t)$，事实上，我们可以把$k(t)$吸收进$\lambda$中，通过增加一项$\int_{0}^{t}k(t')\dif t'$定义一个新的$\lambda$。这样有，
\begin{equation}
    \left.\begin{aligned}
    \boldsymbol{A}'&=\boldsymbol{A}+\boldsymbol{\nabla}\lambda\\
    V'&=V-\frac{\partial\lambda}{\partial t}
    \end{aligned}\right\}
\end{equation}

这种对$V$和$\boldsymbol{A}$的改变称为规范变换。它们可用来调整$\boldsymbol{A}$的散度。在静磁学中，取$\boldsymbol{\nabla}·\boldsymbol{A}=\boldsymbol{0}$是最合适的。在电动力学情形中，
情况就不这么明了，最方便的规范变换依赖于具体的问题。在文献中有许多著名的规范变换。下
面将给出两个最流行的。
\subsubsection{库仑规范与洛伦兹规范}
\mysssec{证明$V$和$\boldsymbol{A}$的微分方程，可以写成下面更对称的形式
\begin{equation}
    \left.\begin{aligned}
    \square^2V+\frac{\partial L}{\partial t}&=-\frac{1}{\varepsilon_0}\rho\\
        \square^2\boldsymbol{A}+\boldsymbol{\nabla}L&=-\mu_0\boldsymbol{J}\\
    \end{aligned}\right\}
\end{equation}
式中
\begin{equation}
    \square^2\equiv\nabla^2-\mu_0\varepsilon_0\frac{\partial^2}{\partial t^2},\,L\equiv\boldsymbol{\nabla}\cdot\boldsymbol{A}+\mu_0\varepsilon_0\frac{\partial V}{\partial t}
\end{equation}
}
\begin{align*}
    -\nabla^2V-\frac{\partial\boldsymbol{\boldsymbol{\nabla}\cdot A}}{\partial t}&=\frac{1}{\varepsilon_0}\rho\\
    \nabla^2V+\frac{\partial \boldsymbol{\nabla}\cdot\boldsymbol{A}}{\partial t}&=-\frac{1}{\varepsilon_0}\rho\\
    \left(\nabla^2-\mu_0\varepsilon_0\frac{\partial^2}{\partial t^2}\right)V+\frac{\partial \boldsymbol{\nabla}\cdot\boldsymbol{A}+\mu_0\varepsilon_0\frac{\partial V}{\partial t}}{\partial t}&=-\frac{1}{\varepsilon_0}\rho\\
    \square^2V+\frac{\partial L}{\partial t}&=-\frac{1}{\varepsilon_0}\rho\\
    \boldsymbol{\nabla}\times\left(\boldsymbol{\nabla}\times\boldsymbol{A}\right)&=\mu_0\boldsymbol{J}
    -\mu_0\varepsilon_0\boldsymbol{\nabla}\frac{\partial V}{\partial t}
    -\mu_0\varepsilon_0\frac{\partial^2\boldsymbol{A}}{\partial t^2}\\
    \boldsymbol{\nabla}\left(\boldsymbol{\nabla}\cdot\boldsymbol{A}\right)
    -\nabla^2\boldsymbol{A}&=\mu_0\boldsymbol{J}
    -\mu_0\varepsilon_0\boldsymbol{\nabla}\frac{\partial V}{\partial t}
    -\mu_0\varepsilon_0\frac{\partial^2\boldsymbol{A}}{\partial t^2}\\
    \left(\nabla^2-\mu_0\varepsilon_0\frac{\partial^2}{\partial t^2}\right)\boldsymbol{A}+\boldsymbol{\nabla}\left(\boldsymbol{\nabla}\cdot\boldsymbol{A}+\mu_0\varepsilon_0\frac{\partial V}{\partial t}\right)&=-\mu_0\boldsymbol{J}\\
    \square^2\boldsymbol{A}+\boldsymbol{\nabla}L&=-\mu_0\boldsymbol{J}
\end{align*}
\mysssec{
    $$
        V=0,\,\boldsymbol{A}=\left\{\begin{aligned}
            &\frac{\mu_0k}{4c}\left(ct-|x|\right)^2\boldsymbol{e}_z,\,&|x|<ct\\
            &0,\,&|x|>ct
        \end{aligned}\right.
    $$
    考虑一边长为$l$，宽为$w$，高为$h$的长方形盒子放在距$yz$面高为$d$
的地方。\\
(a)求在时间$t_1=d/c$和$t_2=(d+h)/c$时，盒子中的能量。\\
(b)求出坡印廷矢量($x>0$)。确定在时间间隔$t_1<t<t_2$内，单位时间流入盒子的能量。}
(a)\begin{align*}
    W&=\frac{1}{2}\varepsilon_0E^2+\frac{1}{2\mu_0}B^2\\
    W&=\frac{1}{2}\varepsilon_0\left(-\boldsymbol{\nabla}V-\frac{\partial\boldsymbol{A}}{\partial t}\right)^2+\frac{1}{2\mu_0}\left(\boldsymbol{\nabla}\times\boldsymbol{A}\right)^2\\
    W&=\frac{1}{2}\varepsilon_0\left(\frac{\partial\boldsymbol{A}}{\partial t}\right)^2+\frac{1}{2\mu_0}\left(\boldsymbol{\nabla}\times\boldsymbol{A}\right)^2\\
    W_1&=0\\
    W_2&=wl\int_{d}^{d+h}\left\{\frac{1}{2}\varepsilon_0\left[\frac{\partial\frac{\mu_0k}{4c}\left(ct-|x|\right)^2\boldsymbol{e}_z}{\partial t}\right]^2+\frac{1}{2\mu_0}\left[\boldsymbol{\nabla}\times\frac{\mu_0k}{4c}\left(ct-|x|\right)^2\boldsymbol{e}_z\right]^2\right\}\dif x\\
    W_2&=wl\int_{d}^{d+h}\left\{\frac{1}{2}\varepsilon_0\left[\frac{\partial\frac{\mu_0k}{4c}\left(ct-x\right)^2}{\partial t}\right]^2+\frac{1}{2\mu_0}\left[\boldsymbol{\nabla}\times\frac{\mu_0k}{4c}\left(d+h-x\right)^2\boldsymbol{e}_z\right]^2\right\}\dif x\\
    W_2&=wl\int_{d}^{d+h}\left\{\frac{1}{2}\varepsilon_0\left[\frac{\mu_0k}{2}\left(ct-x\right)\right]^2+\frac{1}{2\mu_0}\left[\frac{\mu_0k}{2c}\left(x-d-h\right)\right]^2\right\}\dif x\\
    W_2&=wl\int_{d}^{d+h}\left(\varepsilon_0\frac{\mu_0^2k^2}{8}
    +\frac{\mu_0k^2}{8c^2}\right)\left(x-d-h\right)^2\dif x\\
    W_2&=wl\left.\frac{k^2}{12}\frac{\mu_0}{c^2}\left(x-d-h\right)^3\right|_{d}^{d+h}\\
    W_2&=\frac{k^2}{12}\frac{\mu_0}{c^2}wlh^3
\end{align*}

(b)\begin{align*}
    \boldsymbol{S}&=\frac{1}{\mu_0}\boldsymbol{E}\times\boldsymbol{B}\\
    \boldsymbol{S}&=\frac{1}{\mu_0}\left[-\frac{\mu_0k}{2}\left(ct-x\right)\boldsymbol{e}_z\right]\times\left[\frac{\mu_0k}{2c}(ct-x)\boldsymbol{e}_y\right]\\
    \boldsymbol{S}&=\frac{\mu_0k^2}{4c}\left(ct-x\right)^2\boldsymbol{e}_x\\
    \int_{t_1}^{t_2}\boldsymbol{S}\dif t&=wl\int_{t_1}^{t_2}\frac{\mu_0k^2}{4c}\left(ct-d\right)^2\boldsymbol{e}_x\dif t\\
    \int_{t_1}^{t_2}\boldsymbol{S}\dif t&=wl\left.\frac{\mu_0k^2}{12c^2}\left(ct-d\right)^3\boldsymbol{e}_x\right|_{t_1}^{t_2}\\
    \int_{t_1}^{t_2}\boldsymbol{S}\dif t&=\frac{k^2}{12}\frac{\mu_0}{c^2}wlh^3
\end{align*}
\mysssec{$V=0,\,\boldsymbol{A}=-\dfrac{qt}{4\pi\varepsilon_0r^2}\boldsymbol{e}_r$求场分布,\, 利用规范函数$\lambda=-\dfrac{qt}{4\pi\varepsilon_0r}$变换势}
\begin{align*}
    \boldsymbol{E}&=-\boldsymbol{\nabla}V-\frac{\partial\boldsymbol{A}}{\partial t}\\
    \boldsymbol{E}&=\frac{\partial\dfrac{qt}{4\pi\varepsilon_0r^2}\boldsymbol{e}_r}{\partial t}\\
    \boldsymbol{E}&=\dfrac{q}{4\pi\varepsilon_0r^2}\boldsymbol{e}_r\\
    \boldsymbol{B}&=\boldsymbol{\nabla}\times\boldsymbol{A}\\
    \boldsymbol{B}&=-\boldsymbol{\nabla}\times\dfrac{qt}{4\pi\varepsilon_0r^2}\boldsymbol{e}_r\\
    \boldsymbol{B}&=\boldsymbol{0}\\
    \end{align*}
    \begin{align*}
    \boldsymbol{A}'&=\boldsymbol{A}+\boldsymbol{\nabla}\lambda\\
    \boldsymbol{A}'&=-\dfrac{qt}{4\pi\varepsilon_0r^2}\boldsymbol{e}_r-\boldsymbol{\nabla}\dfrac{qt}{4\pi\varepsilon_0r}\\
    \boldsymbol{A}'&=-\dfrac{qt}{4\pi\varepsilon_0r^2}\boldsymbol{e}_r-\dfrac{\partial \dfrac{qt}{4\pi\varepsilon_0r}}{\partial r}\boldsymbol{e}_r\\
    \boldsymbol{A}'&=-\dfrac{qt}{4\pi\varepsilon_0r^2}\boldsymbol{e}_r+\dfrac{qt}{4\pi\varepsilon_0r^2}\boldsymbol{e}_r\\
    \boldsymbol{A}'&=0\\
    V'&=V-\frac{\partial\lambda}{\partial t}\\
    V'&=\frac{\partial\dfrac{qt}{4\pi\varepsilon_0r}}{\partial t}\\
    V'&=\dfrac{q}{4\pi\varepsilon_0r}\\
\end{align*}
\mysssec{设$V=0$，$A=A_0\sin(kx-\omega t)\boldsymbol{e}_y$，其中$A_0$，$\omega$和$k$为常数。求出$\boldsymbol{E}$和$\boldsymbol{B}$，并验证它们满足真
空中的麦克斯韦方程。对$\omega$和$k$必须施加什么条件？}
\begin{align*}
    \boldsymbol{E}&=-\boldsymbol{\nabla}V-\frac{\partial\boldsymbol{A}}{\partial t}\\
    \boldsymbol{E}&=-\frac{\partial A_0\sin(kx-\omega t)\boldsymbol{e}_y}{\partial t}\\
    \boldsymbol{E}&=A_0\omega\cos(kx-\omega t)\boldsymbol{e}_y\\
    \boldsymbol{B}&=\boldsymbol{\nabla}\times\boldsymbol{A}\\
    \boldsymbol{B}&=\boldsymbol{\nabla}\times A_0\sin(kx-\omega t)\boldsymbol{e}_y\\
    \boldsymbol{B}&=A_0k\cos(kx-\omega t)\boldsymbol{e}_z\\
    \end{align*}
    \begin{align*}
    \boldsymbol{\nabla}\cdot\boldsymbol{E}&=\boldsymbol{\nabla}\cdot A_0\omega\cos(kx-\omega t)\boldsymbol{e}_y\\
    \boldsymbol{\nabla}\cdot\boldsymbol{E}&=0\\
    \boldsymbol{\nabla}\times\boldsymbol{E}&=-\frac{\partial \boldsymbol{B}}{\partial t}\\
    \boldsymbol{\nabla}\cdot\boldsymbol{B}&=0\\
    \boldsymbol{\nabla}\times\boldsymbol{B}&=\mu_0\boldsymbol{J}+\mu_0\varepsilon_0\frac{\partial \boldsymbol{E}}{\partial t}\\
    \boldsymbol{\nabla}\times A_0k\cos(kx-\omega t)\boldsymbol{e}_z&=\mu_0\boldsymbol{J}+\mu_0\varepsilon_0\frac{\partial A_0\omega\cos(kx-\omega t)\boldsymbol{e}_y}{\partial t}\\
    A_0k^2\sin(kx-\omega t)\boldsymbol{e}_y&=\mu_0\boldsymbol{J}+\mu_0\varepsilon_0A_0\omega^2\sin(kx-\omega t)\boldsymbol{e}_y\\
    k^2&=\mu_0\varepsilon_0\omega^2
\end{align*}
